% Môn: Vật lý
% Lớp: 10
% Chương: Chương I. Mở đầu
% Bài: L10C1 Bài 2. Các quy tắc an toàn trong phòng thực hành Vật lí · An toàn với nguồn nhiệt và nguy cơ cháy nổ
% Số câu: 185
% App đọc file này trực tiếp — sửa rồi Commit trên GitHub, bấm Đồng bộ GitHub .tex

% L10C1 Bài 2. Các quy tắc an toàn trong phòng thực hành Vật lí



% ===== Câu 39 =====
% ID: L10C1B2-12-DS
% Mức: NB
\dangbt{An toàn với nguồn nhiệt và nguy cơ cháy nổ}
\begin{ex}
% ID: L10C1B2-12-DS
% Mức: NB
Trong tình huống sau khi hoàn thành phép đo đối với nguồn nhiệt và chất dễ cháy, hãy đánh giá các phát biểu sau.
\choiceTF
{\True Cần đặt nguồn nhiệt trên mặt phẳng chịu nhiệt, cách xa chất dễ cháy và chuẩn bị phương tiện dập lửa phù hợp.}
{Có thể đun bình kín hoặc để cồn gần ngọn lửa nếu thao tác diễn ra nhanh.}
{\True Phải báo cho giáo viên hoặc người phụ trách khi phát hiện nguy cơ vượt khả năng kiểm soát.}
{\True Chỉ tiếp tục thí nghiệm sau khi nguyên nhân nguy hiểm đã được loại bỏ hoặc kiểm soát.}
\loigiai{
\begin{itemchoice}
\item (Đúng) Cần đặt nguồn nhiệt trên mặt phẳng chịu nhiệt, cách xa chất dễ cháy và chuẩn bị phương tiện dập lửa phù hợp. (Vì): ây là biện pháp an toàn bắt buộc khi làm việc với nguồn nhiệt và chất dễ cháy: đặt nguồn nhiệt trên mặt phẳng chịu nhiệt, giữ khoảng cách an toàn với chất dễ cháy và luôn chuẩn bị sẵn phương tiện dập lửa phù hợp
\item (Sai) Có thể đun bình kín hoặc để cồn gần ngọn lửa nếu thao tác diễn ra nhanh. (Vì): Đun bình kín có thể gây nổ do áp suất tăng cao; để cồn gần ngọn lửa có thể gây cháy lan. Những hành vi này luôn nguy hiểm dù thao tác diễn ra nhanh hay chậm
\item (Đúng) Phải báo cho giáo viên hoặc người phụ trách khi phát hiện nguy cơ vượt khả năng kiểm soát. (Vì): Khi phát hiện nguy cơ vượt khả năng kiểm soát, bắt buộc phải báo ngay cho giáo viên hoặc người phụ trách để xử lý kịp thời, tránh để sự cố lan rộng
\item (Đúng) Chỉ tiếp tục thí nghiệm sau khi nguyên nhân nguy hiểm đã được loại bỏ hoặc kiểm soát. (Vì): Chỉ được tiếp tục thí nghiệm sau khi nguyên nhân gây nguy hiểm đã được loại bỏ hoặc kiểm soát hoàn toàn, đảm bảo môi trường thí nghiệm an toàn
\end{itemchoice}
}
\end{ex}

% ===== Câu 50 =====
% ID: L10C1B2-13-DS
% Mức: TH
\begin{ex}
% ID: L10C1B2-13-DS
% Mức: TH
Trong tình huống khi một bạn chưa đọc hướng dẫn đối với nguồn nhiệt và chất dễ cháy, hãy đánh giá các phát biểu sau.
\choiceTF
{\True Cần đặt nguồn nhiệt trên mặt phẳng chịu nhiệt, cách xa chất dễ cháy và chuẩn bị phương tiện dập lửa phù hợp.}
{Có thể đun bình kín hoặc để cồn gần ngọn lửa nếu thao tác diễn ra nhanh.}
{\True Phải báo cho giáo viên hoặc người phụ trách khi phát hiện nguy cơ vượt khả năng kiểm soát.}
{\True Chỉ tiếp tục thí nghiệm sau khi nguyên nhân nguy hiểm đã được loại bỏ hoặc kiểm soát.}
\loigiai{
\begin{itemchoice}
\item (Đúng) Cần đặt nguồn nhiệt trên mặt phẳng chịu nhiệt, cách xa chất dễ cháy và chuẩn bị phương tiện dập lửa phù hợp. (Vì): vì đây là biện pháp an toàn của dạng nguồn nhiệt và chất dễ cháy; b sai vì hành vi “đun bình kín hoặc để cồn gần ngọn lửa” vẫn nguy hiểm; c và d đúng do sự cố phải được báo cáo, cô lập và kiểm soát trước khi tiếp tục
\item (Sai) Có thể đun bình kín hoặc để cồn gần ngọn lửa nếu thao tác diễn ra nhanh.
\item (Đúng) Phải báo cho giáo viên hoặc người phụ trách khi phát hiện nguy cơ vượt khả năng kiểm soát.
\item (Đúng) Chỉ tiếp tục thí nghiệm sau khi nguyên nhân nguy hiểm đã được loại bỏ hoặc kiểm soát.
\end{itemchoice}
}
\end{ex}

% ===== Câu 51 =====
% ID: L10C1B2-13-TL
% Mức: TH
\begin{bt}
% ID: L10C1B2-13-TL
% Mức: TH
Trình bày nguyên tắc cốt lõi khi làm việc với nguồn nhiệt và chất dễ cháy và giải thích vì sao phải tuân thủ.
\loigiai{
Nguyên tắc cốt lõi là đặt nguồn nhiệt trên mặt phẳng chịu nhiệt, cách xa chất dễ cháy và chuẩn bị phương tiện dập lửa phù hợp. Phải tuân thủ vì tránh bỏng, tăng áp suất gây nổ và cháy lan; đồng thời luôn dừng và báo người phụ trách khi điều kiện không bảo đảm.
}
\end{bt}

% ===== Câu 52 =====
% ID: L10C1B2-13-TLN
% Mức: NB
\begin{ex}
% ID: L10C1B2-13-TLN
% Mức: NB
Một quy trình về nguồn nhiệt và chất dễ cháy gồm: (1) nhận diện nguy cơ; (2) đặt nguồn nhiệt trên mặt phẳng chịu nhiệt, cách xa chất dễ cháy và chuẩn bị phương tiện dập lửa phù hợp; (3) đun bình kín hoặc để cồn gần ngọn lửa; (4) báo giáo viên khi có bất thường; (5) chỉ tiếp tục khi nguy cơ đã được kiểm soát. Có bao nhiêu bước phù hợp quy tắc an toàn? Chỉ ghi số.
\shortans{4}
\loigiai{
Các bước (1), (2), (4), (5) phù hợp; bước (3) không an toàn. Có 4 bước phù hợp.
}
\end{ex}

% ===== Câu 63 =====
% ID: L10C1B2-14-DS
% Mức: TH
\begin{ex}
% ID: L10C1B2-14-DS
% Mức: TH
Trong tình huống khi khu vực làm việc bị ướt đối với nguồn nhiệt và chất dễ cháy, hãy đánh giá các phát biểu sau.
\choiceTF
{\True Cần đặt nguồn nhiệt trên mặt phẳng chịu nhiệt, cách xa chất dễ cháy và chuẩn bị phương tiện dập lửa phù hợp.}
{Có thể đun bình kín hoặc để cồn gần ngọn lửa nếu thao tác diễn ra nhanh.}
{\True Phải báo cho giáo viên hoặc người phụ trách khi phát hiện nguy cơ vượt khả năng kiểm soát.}
{\True Chỉ tiếp tục thí nghiệm sau khi nguyên nhân nguy hiểm đã được loại bỏ hoặc kiểm soát.}
\loigiai{
\begin{itemchoice}
\item (Đúng) Cần đặt nguồn nhiệt trên mặt phẳng chịu nhiệt, cách xa chất dễ cháy và chuẩn bị phương tiện dập lửa phù hợp. (Vì): vì đây là biện pháp an toàn của dạng nguồn nhiệt và chất dễ cháy; b sai vì hành vi “đun bình kín hoặc để cồn gần ngọn lửa” vẫn nguy hiểm; c và d đúng do sự cố phải được báo cáo, cô lập và kiểm soát trước khi tiếp tục
\item (Sai) Có thể đun bình kín hoặc để cồn gần ngọn lửa nếu thao tác diễn ra nhanh.
\item (Đúng) Phải báo cho giáo viên hoặc người phụ trách khi phát hiện nguy cơ vượt khả năng kiểm soát.
\item (Đúng) Chỉ tiếp tục thí nghiệm sau khi nguyên nhân nguy hiểm đã được loại bỏ hoặc kiểm soát.
\end{itemchoice}
}
\end{ex}

% ===== Câu 64 =====
% ID: L10C1B2-14-TL
% Mức: TH
\begin{bt}
% ID: L10C1B2-14-TL
% Mức: TH
Phân tích hành vi “đun bình kín hoặc để cồn gần ngọn lửa” và đề xuất cách thay thế an toàn.
\loigiai{
Hành vi “đun bình kín hoặc để cồn gần ngọn lửa” có thể gây tai nạn vì làm mất kiểm soát nguồn nguy hiểm. Cần dừng thao tác, cô lập nguy cơ và thay bằng biện pháp: đặt nguồn nhiệt trên mặt phẳng chịu nhiệt, cách xa chất dễ cháy và chuẩn bị phương tiện dập lửa phù hợp.
}
\end{bt}

% ===== Câu 73 =====
% ID: L10C1B2-15-DS
% Mức: VD
\begin{ex}
% ID: L10C1B2-15-DS
% Mức: VD
Trong tình huống khi cần rời bàn thí nghiệm đối với nguồn nhiệt và chất dễ cháy, hãy đánh giá các phát biểu sau.
\choiceTF
{\True Cần đặt nguồn nhiệt trên mặt phẳng chịu nhiệt, cách xa chất dễ cháy và chuẩn bị phương tiện dập lửa phù hợp.}
{Có thể đun bình kín hoặc để cồn gần ngọn lửa nếu thao tác diễn ra nhanh.}
{\True Phải báo cho giáo viên hoặc người phụ trách khi phát hiện nguy cơ vượt khả năng kiểm soát.}
{\True Chỉ tiếp tục thí nghiệm sau khi nguyên nhân nguy hiểm đã được loại bỏ hoặc kiểm soát.}
\loigiai{
\begin{itemchoice}
\item (Đúng) Cần đặt nguồn nhiệt trên mặt phẳng chịu nhiệt, cách xa chất dễ cháy và chuẩn bị phương tiện dập lửa phù hợp. (Vì): vì đây là biện pháp an toàn của dạng nguồn nhiệt và chất dễ cháy; b sai vì hành vi “đun bình kín hoặc để cồn gần ngọn lửa” vẫn nguy hiểm; c và d đúng do sự cố phải được báo cáo, cô lập và kiểm soát trước khi tiếp tục
\item (Sai) Có thể đun bình kín hoặc để cồn gần ngọn lửa nếu thao tác diễn ra nhanh.
\item (Đúng) Phải báo cho giáo viên hoặc người phụ trách khi phát hiện nguy cơ vượt khả năng kiểm soát.
\item (Đúng) Chỉ tiếp tục thí nghiệm sau khi nguyên nhân nguy hiểm đã được loại bỏ hoặc kiểm soát.
\end{itemchoice}
}
\end{ex}

% ===== Câu 74 =====
% ID: L10C1B2-15-TL
% Mức: VD
\begin{bt}
% ID: L10C1B2-15-TL
% Mức: VD
Xây dựng một bảng kiểm ngắn gồm ít nhất bốn mục cho tình huống khi một bạn chưa đọc hướng dẫn liên quan đến nguồn nhiệt và chất dễ cháy.
\loigiai{
Bảng kiểm gồm: đọc hướng dẫn; nhận diện mối nguy; kiểm tra dụng cụ và bảo hộ; thực hiện đặt nguồn nhiệt trên mặt phẳng chịu nhiệt, cách xa chất dễ cháy và chuẩn bị phương tiện dập lửa phù hợp; theo dõi bất thường; thu dọn và báo cáo sau thí nghiệm.
}
\end{bt}

% ===== Câu 76 =====
% ID: L10C1B2-16-TL
% Mức: VD
\begin{bt}
% ID: L10C1B2-16-TL
% Mức: VD
Một học sinh đã đun bình kín hoặc để cồn gần ngọn lửa. Hãy nêu trình tự xử lí ngay sau khi phát hiện hành vi này.
\loigiai{
Trình tự: yêu cầu dừng ngay; không tiếp cận nếu chưa an toàn; cô lập nguồn nguy hiểm; báo giáo viên; sơ cứu hoặc xử lí theo hướng dẫn; chỉ hoạt động lại sau khi được xác nhận an toàn.
}
\end{bt}

% ===== Câu 79 =====
% ID: L10C1B2-17-TL
% Mức: VD
\begin{bt}
% ID: L10C1B2-17-TL
% Mức: VD
So sánh biện pháp phòng ngừa và biện pháp ứng phó đối với nguy cơ thuộc nhóm nguồn nhiệt và chất dễ cháy.
\loigiai{
Phòng ngừa diễn ra trước sự cố, tập trung vào kiểm tra, loại bỏ hoặc giảm nguy cơ. Ứng phó diễn ra khi sự cố đã hoặc sắp xảy ra, gồm dừng hoạt động, cô lập nguồn nguy hiểm, báo cáo và sơ cứu. Hai nhóm biện pháp bổ sung cho nhau.
}
\end{bt}

% ===== Câu 82 =====
% ID: L10C1B2-18-TL
% Mức: VDC
\begin{bt}
% ID: L10C1B2-18-TL
% Mức: VDC
Đề xuất cách tổ chức nhóm học sinh để thực hiện nhiệm vụ về nguồn nhiệt và chất dễ cháy an toàn và có thể kiểm tra được.
\loigiai{
Phân công trưởng nhóm kiểm tra điều kiện, người thao tác, người quan sát an toàn và người ghi chép. Dùng bảng kiểm, xác nhận trước mỗi bước, quy định tín hiệu dừng và báo ngay bất thường.
}
\end{bt}

% ===== Câu 92 =====
% ID: L10C1B2-22-TN
% Mức: NB
\begin{ex}
% ID: L10C1B2-22-TN
% Mức: NB
Khi phát hiện thiết bị có dấu hiệu bất thường liên quan đến nguồn nhiệt và chất dễ cháy, hành động nào an toàn nhất?
\choice
{đun bình kín hoặc để cồn gần ngọn lửa}
{tiếp tục thao tác rồi báo cáo sau}
{nhờ bạn khác làm thay mà không kiểm tra điều kiện an toàn}
{\True đặt nguồn nhiệt trên mặt phẳng chịu nhiệt, cách xa chất dễ cháy và chuẩn bị phương tiện dập lửa phù hợp}
\loigiai{
Hành động đúng là “đặt nguồn nhiệt trên mặt phẳng chịu nhiệt, cách xa chất dễ cháy và chuẩn bị phương tiện dập lửa phù hợp” vì tránh bỏng, tăng áp suất gây nổ và cháy lan. Chọn D.
}
\end{ex}

% ===== Câu 94 =====
% ID: L10C1B2-23-TN
% Mức: NB
\begin{ex}
% ID: L10C1B2-23-TN
% Mức: NB
Trong lúc thay đổi cách bố trí dụng cụ liên quan đến nguồn nhiệt và chất dễ cháy, hành động nào an toàn nhất?
\choice
{tiếp tục thao tác rồi báo cáo sau}
{nhờ bạn khác làm thay mà không kiểm tra điều kiện an toàn}
{\True đặt nguồn nhiệt trên mặt phẳng chịu nhiệt, cách xa chất dễ cháy và chuẩn bị phương tiện dập lửa phù hợp}
{đun bình kín hoặc để cồn gần ngọn lửa}
\loigiai{
Hành động đúng là “đặt nguồn nhiệt trên mặt phẳng chịu nhiệt, cách xa chất dễ cháy và chuẩn bị phương tiện dập lửa phù hợp” vì tránh bỏng, tăng áp suất gây nổ và cháy lan. Chọn C.
}
\end{ex}

% ===== Câu 96 =====
% ID: L10C1B2-24-TN
% Mức: TH
\begin{ex}
% ID: L10C1B2-24-TN
% Mức: TH
Sau khi hoàn thành phép đo liên quan đến nguồn nhiệt và chất dễ cháy, hành động nào an toàn nhất?
\choice
{nhờ bạn khác làm thay mà không kiểm tra điều kiện an toàn}
{\True đặt nguồn nhiệt trên mặt phẳng chịu nhiệt, cách xa chất dễ cháy và chuẩn bị phương tiện dập lửa phù hợp}
{đun bình kín hoặc để cồn gần ngọn lửa}
{tiếp tục thao tác rồi báo cáo sau}
\loigiai{
Hành động đúng là “đặt nguồn nhiệt trên mặt phẳng chịu nhiệt, cách xa chất dễ cháy và chuẩn bị phương tiện dập lửa phù hợp” vì tránh bỏng, tăng áp suất gây nổ và cháy lan. Chọn B.
}
\end{ex}

% ===== Câu 98 =====
% ID: L10C1B2-25-TN
% Mức: TH
\begin{ex}
% ID: L10C1B2-25-TN
% Mức: TH
Khi một bạn chưa đọc hướng dẫn liên quan đến nguồn nhiệt và chất dễ cháy, hành động nào an toàn nhất?
\choice
{\True đặt nguồn nhiệt trên mặt phẳng chịu nhiệt, cách xa chất dễ cháy và chuẩn bị phương tiện dập lửa phù hợp}
{đun bình kín hoặc để cồn gần ngọn lửa}
{tiếp tục thao tác rồi báo cáo sau}
{nhờ bạn khác làm thay mà không kiểm tra điều kiện an toàn}
\loigiai{
Hành động đúng là “đặt nguồn nhiệt trên mặt phẳng chịu nhiệt, cách xa chất dễ cháy và chuẩn bị phương tiện dập lửa phù hợp” vì tránh bỏng, tăng áp suất gây nổ và cháy lan. Chọn A.
}
\end{ex}

% ===== Câu 102 =====
% ID: L10C1B2-27-TN
% Mức: TH
\begin{ex}
% ID: L10C1B2-27-TN
% Mức: TH
Khi cần rời bàn thí nghiệm liên quan đến nguồn nhiệt và chất dễ cháy, hành động nào an toàn nhất?
\choice
{tiếp tục thao tác rồi báo cáo sau}
{nhờ bạn khác làm thay mà không kiểm tra điều kiện an toàn}
{\True đặt nguồn nhiệt trên mặt phẳng chịu nhiệt, cách xa chất dễ cháy và chuẩn bị phương tiện dập lửa phù hợp}
{đun bình kín hoặc để cồn gần ngọn lửa}
\loigiai{
Hành động đúng là “đặt nguồn nhiệt trên mặt phẳng chịu nhiệt, cách xa chất dễ cháy và chuẩn bị phương tiện dập lửa phù hợp” vì tránh bỏng, tăng áp suất gây nổ và cháy lan. Chọn C.
}
\end{ex}

% ===== Câu 105 =====
% ID: L10C1B2-29-TN
% Mức: VD
\begin{ex}
% ID: L10C1B2-29-TN
% Mức: VD
Khi kết quả đo khác dự kiến liên quan đến nguồn nhiệt và chất dễ cháy, hành động nào an toàn nhất?
\choice
{\True đặt nguồn nhiệt trên mặt phẳng chịu nhiệt, cách xa chất dễ cháy và chuẩn bị phương tiện dập lửa phù hợp}
{đun bình kín hoặc để cồn gần ngọn lửa}
{tiếp tục thao tác rồi báo cáo sau}
{nhờ bạn khác làm thay mà không kiểm tra điều kiện an toàn}
\loigiai{
Hành động đúng là “đặt nguồn nhiệt trên mặt phẳng chịu nhiệt, cách xa chất dễ cháy và chuẩn bị phương tiện dập lửa phù hợp” vì tránh bỏng, tăng áp suất gây nổ và cháy lan. Chọn A.
}
\end{ex}

